% 
% \eric{general intro to agentic coding, multi-agent coding, etc.}
% \textbf{Agentic Coding Systems using Large Language Models}
% The language model community has established a robust foundation in agentic system development by pairing a central "brain" LLM with executable functions. This "Tool-Use" paradigm has achieved remarkable success in automating single-session tasks such as webpage building, autonomous search, and natural language dialogue .

% - Context window 

% - Multi agent

% - Skill library 

% \textbf{Graph in Robotics} TAMP/(Modular Robotic Planning Systems)
% /ROS


% GaP extends CaP and CaP-X in the following ways:   instead of generating a single Python script, it represents the policy as a typed directed computation graph, where nodes are modular perception, planning, control, or verification operations, and edges encode their dependencies. This structure targets known industrial automation task classes with repeated execution under structured variation, where the goal is to produce a reliable workflow rather than solve only a single trial. Second, GaP uses an explicit multi-agent architecture: a coordinator decomposes the task, skill agents generate local subgraphs, and the coordinator aggregates and verifies them into a complete policy graph. Third, GaP introduces simulation-based verification as a refinement mechanism. It recreates the task class in simulation, evaluates candidate graphs through rollouts, and uses physical success or failure to revise graph connectivity, node parameters, and skill choices before deployment. In this sense, simulation serves as a verification step for physical success: it preemptively checks whether the composed graph produces a reliable workflow under structured variation.

\section{Related Work}

\textbf{Variational Automation}.
The term ``automation'' is commonly used to describe robot systems that persistently perform repetitive tasks over hours, weeks, months, or years for logistics, manufacturing, healthcare, service, agriculture, and other applications. To be successful, commercial and industrial automation must achieve desired throughput (units per hour: success rate divided by cycle time) at a cost that will provide a desired return on investment. Accordingly, research in automation extends research in robotics by focusing on reliability, cost, safety, ease of use, and other factors required for successful production deployment.

In this paper, we propose ``Variational Automation (VA)'' as a class of tasks that differ from Generalist Robotics (GR) and Fixed Automation (FA). In FA, a robot persistently performs identical instances of a task with objects of identical geometry. In box stacking, for example, FA would consist of moving identical boxes from a common initial pose on a conveyor belt into a fixed pallet arrangement. In FA, variations in the environment and object shape and pose are minimal. 

For GR, a robot performs a variety of different tasks (e.g., placing groceries into a refrigerator, or 
domestic cleaning, folding, kitchen pick-and-place) in different homes, where environments vary considerably and objects have variable geometry and highly variable initial poses. Recent robot learning research predominantly targets highly unstructured Generalist Robotics (GR), such as household robots executing diverse chores across novel environments using a single generalist model-free VLA policies~\cite{levine2016end, chi2025diffusion, kim24openvla, pi0_2024, black2025pi_}.

In VA, by contrast, a robot persistently performs varying instances of the task: the boxes have variable geometry (different SKUs), arrive in a distribution of varying initial poses, and must be densely packed into varying pallet arrangements. As most robot learning research focuses on GR, and robot learning may not be required for FA, we focus on robot learning for VA, aiming to reduce the human effort required to set up VA systems.

\textbf{Modular Robot Control Methods}.
Early robotic systems programmed policies using classical modular components. The Shakey robot~\cite{nilsson1998artificial} in the late 1960s
decomposed its architecture into 3 functional components:
sensing, planning, and execution~\cite{nilsson2014principles}. System architectures like the Robot Operating System (ROS) \cite{macenski2022robot} utilize explicit directed graphs to route data, manage dependencies, and ensure reliable execution.
Concurrently, TAMP \cite{kaelbling2011hierarchical,garrett2021integrated} approaches jointly solve discrete task planning and continuous motion planning problems, enabling the satisfaction of constraints involving both high-level action sequencing and low-level geometric feasibility.

% Recently, the rise of Large Language Models has catalyzed a shift toward generative task decomposition. Frameworks such as

Code-as-Policy (CaP) uses agents to write robot control code, suggesting an alternative to manual coding of model-based methods and pure end-to-end learning of model-free policies~\cite{codeaspolicies2022}. 
Extensions such as CaP-X \cite{fu2026cap}, GRAPPA~\cite{bucker2026grappa} and Maestro \cite{shi2025maestro} use more recent LLMs to generate robot code. Similarly, systems like TiPToP~\cite{shen2026tiptop} incorporate LLMs with classical robot Task and Motion Planning (TAMP)~\cite{wang2024llm, curtis2025trust, yang2025guiding, kumar2026open}. 



% Recently, the rise of Large Language Models has catalyzed a shift toward generative task decomposition. Frameworks such as
Code-as-Policy (CaP)
is an alternative to manual coding of model-based methods and pure end-to-end learning of model-free policies \cite{codeaspolicies2022}, ALGARA \cite{coteagentic}, and CodeDiffuser~\cite{yin2025codediffuser} utilized LLMs for open-vocabulary script generation to generate robot code.
%*******  WHO BUTCHERED this section?  
GaP tempers the flexibility of CaP approaches by introducing a graph structure, allowing GaP to harness the open-world adaptivity of pretrained LLM coding agents while maintaining a structured, interpretable graph structure to support persistent Variational Automation.  
% By bridging the foundational graph representations of systems like ROS with generative agents, we introduce an explicit Graph-as-Policy representation 
% More recent architectures, including CaP-X \cite{fu2026cap} and Maestro \cite{shi2025maestro}, expand upon this foundation by curating extensive, modular skill libraries that range from vision-language perception to motion planning. 

\textbf{Self-Improving Agentic Workflows For Robotic Control}.
Outside of robotics, Large Language Models (LLMs) have demonstrated exceptional capabilities in dynamic coding, complex software integration, and API orchestration \cite{wu2024introducing, team2026qwen3, hou2025model, openai2024gpt4o, anthropic2024claude35, google2025gemini25}. By embedding these capabilities into structured agentic workflows, LLMs can autonomously synthesize solutions for open-ended tasks~\cite{zhang2025aflow}. A foundational example is Voyager \cite{wang2023voyager}, which uses LLMs to iteratively write, refine, and execute Minecraft game playing skills to continually expand a curated library of reusable skills. In these frameworks, a ``skill'' is typically defined as a discrete function with strict semantic contracts. Recent literature has begun to explore the iterative optimization of agentic systems through language-based failure reasoning \cite{zehle2026promptolution} and multi-agent prompt optimization frameworks \cite{agrawal2025gepa, lee2025feedback}.
% In the domain of robotics, an agentic tool-use paradigm translates into a dynamic sequencing and parameterization of model-based or model-free sensorimotor primitives. Early works successfully grounded LLM reasoning in robotic affordances \cite{ahn2022can, huang2022inner}, while frameworks such as 


To improve code generation quality and performance, CaP-X~\cite{fu2026cap} used a VLM to provide feedback before and after execution (i.e. Visual Differencing), but the VLM can suffer from hallucinations and cannot handle geometric and numerical information such as motion feasibility. 
Building on the structure of Blox-Net~\cite{goldberg2025bloxnet}, which uses physical experiments to improve LLM-generated task plans, GaP uses multiple LLM agents to generate a robot computation graph and iteratively improve it using simulation experiments.
% However, while these frameworks generate plausible open-loop plans, they often lack the physical grounding necessary to verify geometric and kinematic feasibility before real-world deployment. To address this reality gap, 
